\section{Model Compression}

\subsection{Tổng quan Model Compression}

Model Compression là nhóm kỹ thuật giảm chi phí lưu trữ hoặc suy luận của mạng
nơ-ron. Các hướng phổ biến gồm pruning, quantization, knowledge distillation và
thiết kế kiến trúc gọn. Đồ án tập trung vào pruning và quantization vì hai kỹ thuật
này có thể áp dụng trực tiếp lên baseline và kết hợp với TensorRT.

Hiệu quả nén phải được đánh giá trên cả kích thước artifact, số phép tính lý
thuyết và latency thực tế. Một mô hình có nhiều trọng số bằng không chưa chắc
chạy nhanh hơn nếu backend vẫn thực hiện phép tính dense.

\subsection{Model Pruning}

Pruning loại bỏ trọng số, kênh hoặc cấu trúc được xem là ít quan trọng. Tiêu chí
có thể dựa trên độ lớn trọng số, chuẩn của filter, độ nhạy loss hoặc thông tin
gradient. Quy trình cơ bản gồm huấn luyện baseline, tính điểm quan trọng, cắt tỉa,
fine-tuning và đánh giá lại \cite{han2016,molchanov2017}.

\subsection{Structured Pruning}

Structured pruning loại bỏ đơn vị có cấu trúc như channel, filter hoặc block.
Kết quả là tensor có kích thước nhỏ hơn và có thể được thư viện dense hiện có
khai thác mà không cần kernel thưa chuyên dụng. Đây là lựa chọn phù hợp khi mục
tiêu cuối là TensorRT trên Jetson.

Nhược điểm là mỗi quyết định loại bỏ channel ảnh hưởng đến các tầng lân cận và
có thể làm accuracy giảm nhanh. Việc cắt tỉa cần tôn trọng quan hệ residual,
concatenation và detection head trong kiến trúc YOLO.

\subsection{Unstructured Pruning}

Unstructured pruning đặt các trọng số riêng lẻ về không và có thể đạt sparsity
cao với mức suy giảm accuracy thấp hơn. Tuy nhiên, kích thước tensor không đổi;
latency chỉ giảm rõ khi phần cứng và runtime hỗ trợ sparse kernel tương ứng.

Trong đồ án, unstructured pruning có thể được dùng làm đối chứng về sparsity,
nhưng cấu hình triển khai ưu tiên structured pruning nếu mục tiêu là giảm latency
thực tế trên Jetson.

\subsection{Pruning Ratio}

Pruning ratio biểu diễn tỷ lệ phần tử hoặc cấu trúc bị loại bỏ:
\begin{equation}
  r=\frac{N_{before}-N_{after}}{N_{before}}\times 100\%.
\end{equation}
Các mức 20\%, 30\% và 40\% được dùng làm điểm thực nghiệm ban đầu. Tỷ lệ này
không mặc định áp dụng đồng đều cho mọi tầng; các tầng đầu, tầng đầu ra và nhánh
nhạy cảm có thể được bảo vệ hoặc đặt ngưỡng riêng.

\subsection{Fine-tuning sau Pruning}

Sau pruning, phân bố đặc trưng bị thay đổi nên mô hình cần được fine-tune với
learning rate nhỏ hơn baseline. Fine-tuning phục hồi accuracy trong khi duy trì
cấu trúc đã cắt. Checkpoint tốt nhất được chọn theo validation mAP và phải được
đánh giá lại trên tập test trước khi chuyển sang QAT hoặc TensorRT.